% Created 2026-04-04 Sat 14:45
% Intended LaTeX compiler: xelatex
\documentclass[11pt,letterpaper]{article}
\usepackage{graphicx}
\usepackage{longtable}
\usepackage{wrapfig}
\usepackage{rotating}
\usepackage[normalem]{ulem}
\usepackage{capt-of}
\usepackage{hyperref}
\usepackage{amsmath}
\usepackage{amssymb}
\usepackage{geometry}
\geometry{left=1.2in, right=1.2in, top=1.0in, bottom=1.0in}
\usepackage{fancyhdr}
\usepackage{booktabs}
\usepackage{array}
\usepackage{setspace}
\usepackage{parskip}
\usepackage{microtype}
\usepackage[hidelinks]{hyperref}
\setlength{\parindent}{0pt}
\setlength{\parskip}{6pt}
% --- Page style: no header on title/TOC, header from page 3 onward ---
\fancypagestyle{plain}{\fancyhf{}\fancyfoot[C]{\thepage}\renewcommand{\headrulewidth}{0pt}}
\fancypagestyle{body}{%
\fancyhf{}%
\fancyhead[C]{\small Version 0.4 $\cdot$ DOI: 10.5281/zenodo.######## $\cdot$ CC BY 4.0}%
\fancyfoot[C]{\thepage}%
\renewcommand{\headrulewidth}{0.4pt}}
\author{Bradley K. DePuy, Independent Researcher}
\date{2026}
\title{V01 The Medium: A Foundational Theory of Reality}
\hypersetup{
 pdfauthor={Bradley K. DePuy, Independent Researcher},
 pdftitle={V01 The Medium: A Foundational Theory of Reality},
 pdfkeywords={},
 pdfsubject={},
 pdfcreator={Emacs 29.3 (Org mode 9.7.39)}, 
 pdflang={English}}
\begin{document}


%% ---- TITLE PAGE (page 1, no header) ----
\pagestyle{plain}
\begin{titlepage}
\vspace*{4cm}
\begin{center}
  {\LARGE\bfseries The Medium}\\[0.6em]
  {\large A Foundational Theory of Reality}\\[2.5em]
  {\normalsize Bradley K. DePuy, Independent Researcher}\\[0.5em]
  {\normalsize ORCID: \mbox{\texttt{0009-0001-0328-2299}}}\\[0.5em]
  {\normalsize 2026}
\end{center}
\end{titlepage}

%% ---- Body begins, header active ----
\pagestyle{body}

The Medium
A New Starting Point
\section{Philosophical Introduction}
\label{sec:org963d72a}
\subsection{The Wall}
\label{sec:org9dc8f35}
Every foundational framework in physics confronts a wall — a point beyond which derivation cannot reach and explanation cannot go. This is not a failure of any particular theory. It is the universal condition of formal reasoning, demonstrated by Gödel and confirmed by the history of every serious attempt to reach the ground of physical existence. The wall is always there. What differs between frameworks is not whether they hit it, but whether they acknowledge it. The Medium begins at the wall — not to climb it, but to locate it precisely, state it honestly, and build everything above it with the rigor that honesty makes possible.
Newton's absolute space and time were not derived — they were assumed. The quantum mechanical Born rule, which connects the wave function to observable probabilities, has no derivation from deeper principles. The symmetry groups of the Standard Model — SU(3) × SU(2) × U(1) — were chosen because they fit the data, not derived from something more primitive. Loop quantum gravity assumes spacetime exists before asking how it is quantized. Every framework, however sophisticated its mathematics, rests on something it cannot account for. The wall is not inside any particular theory. It is outside all of them simultaneously.
Gödel proved this formally in 1931. Any sufficiently powerful formal system contains true statements it cannot prove from within itself. The foundation cannot be self-validating. Something always escapes. This is not a contingent limitation — it is structural. The wall is built into the architecture of formal reasoning itself, not into any particular application of it.
\subsection{The Wall Does Not Move}
\label{sec:org56470f5}
Changing the observational perspective does not move the wall. It is always there, regardless of where one stands or what one assumes. What changes with a change of position is not the wall but the terrain — the problems visible from that position, the paths available for approach, the solutions that become conceivable.
A mountain has a sheer face. Climbers have attempted that face for generations without success. The wall is real. But the mountain has other faces — not because anyone moved the mountain, but because they changed where they were standing. From a different position, routes become visible that were literally out of sight from the original approach. The summit remains where it is. What changes is the terrain available for the attempt.
This is the contribution a different starting point can make. Not a solution to the wall. A different position from which different paths become visible — paths that may approach unsolved problems from angles that existing frameworks, by virtue of their starting positions, cannot access. Whether those paths lead further than existing ones is an open question. The contribution is the position itself, honestly stated, carefully derived, with its own wall precisely located.
\subsection{The Observational Requirement}
\label{sec:org5844b95}
Every current foundational framework is built to explain what is observed. Observation is the court of final appeal. A theory succeeds if it predicts what we see and fails if it does not. This requirement seems obviously correct. It is the definition of science.
But at the foundational level, this requirement carries a hidden circularity. Everything we observe is filtered through the structure of spacetime, the structure of quantum mechanics, and the structure of measurement apparatus — the very structures we are trying to explain. When we demand that a foundational framework explain observations, we are demanding that the foundation match the downstream. We are working backward from consequences and calling that derivation.
The alternative is not to abandon contact with observation. It is to change the relationship between framework and observation. A foundational framework is not validated by explaining what is observed. It is validated by showing that what is observed is among the necessary consequences of what must exist. Observation becomes a consistency check rather than a requirement. The framework is not built to fit the data. It is built from first principles and then compared to the data.
This is a different kind of scientific contribution — and an honest one. It does not claim to explain the universe from scratch. It claims to show that a specific primitive, followed rigorously, produces structures consistent with what we observe. That is achievable. It is also, at the foundational level, the most that is achievable.


Bradley K. DePuy\\
ORCID: 0009-0001-0328-2299\\
Zenodo DOI: 10.5281/zenodo.\#\#\#\#\#\#\#\#
\section{Why This Ontological Approach}
\label{sec:org3535876}

Integrating the Observer‑Bound Bias

The foundational question in cosmology has traditionally been framed as “How did the universe begin?” Yet this question presupposes the very structures we aim to explain: time, causality, geometry, and initial conditions. These features belong to the emergent universe as experienced by observers embedded within it. They are not guaranteed to reflect the universe’s fundamental ontology. By asking how the universe began, we inadvertently assume that time is primitive, that causal chains extend indefinitely backward, and that a “first moment” is a meaningful concept. This framing imports the biases of our observational standpoint into the foundations of theory.

This observer‑bound bias is subtle but pervasive. Every model in modern physics is constructed from the vantage point of observers who exist inside spacetime, using instruments that operate within spacetime, and interpreting data through conceptual frameworks shaped by biological, cognitive, and technological constraints. As a result, physics has accumulated an elaborate superstructure of explanatory devices—singularities, inflationary epochs, vacuum states, symmetry breakings, finely tuned parameters—designed to reconcile emergent observations with models that were never meant to describe the substrate itself. These constructs succeed phenomenologically, but they do so by multiplying assumptions. They reflect the complexity of the world as seen from within, not the simplicity that may characterize the world at its base.

Occam’s razor challenges this trajectory. If the observable universe can arise from a simpler generative substrate, then the additional structures—however mathematically effective—are not fundamental. They are artifacts of describing the universe from the inside. The more basic question is not how the universe began, but what the universe is when stripped of the emergent scaffolding that shapes our perception.

This motivates an ontological approach that begins with the minimal possible primitive: actualization itself. Nothing else is assumed. No space, no time, no geometry, no fields, no laws, no symmetries. Only the bare possibility that something can occur. From this single primitive, the universe is understood not as a collection of objects but as a generative process—a continual unfolding of discrete actualization events whose large‑scale statistical patterns give rise to the structures we observe.

In this framework, spacetime is emergent rather than presupposed. Particles and fields are stable relational patterns rather than fundamental entities. Laws are regularities of co‑manifestation rather than metaphysical absolutes. Even the notion of a “beginning” becomes an emergent feature of the large‑scale pattern, not a primitive fact requiring explanation.

This ontological stance is adopted because it offers the simplest generative substrate capable of producing a universe like ours. It replaces the historical fixation on origins with a deeper inquiry into the nature of the substrate from which origins, time, and structure emerge. By reframing the problem in this way, we avoid the distortions imposed by our embedded perspective and open the possibility of explaining the universe’s richness from a foundation of maximal parsimony.
\section{The Eternal Ground}
\label{sec:orge7a9b2d}
The first medium always was, always is, and always will be.
This formulation uses temporal language to point at something prior to time. It should be understood as pointing rather than defining. More precisely: the medium is not a member of the class of things that begin or end. It is the condition for that class. It has no initial state to explain, no prior cause to identify, no moment zero requiring justification.
Absolute nothingness is not a stable state. Stability requires a constraint — a prohibition against structure arising. But a constraint is already something. Therefore absolute nothing cannot obtain. Therefore something obtains necessarily. That something has no preferred configuration, no structure, no bias — because structure requires prior justification, and there is none available. Strip away everything that requires prior justification and what remains is pure unconstrained indeterminacy. That is the stochastic ground. Not by assumption — by elimination of everything else.
Aristotle's unmoved mover was needed to answer: what started it? If nothing started it — if the generative ground simply is, prior to the conditions that make starting meaningful — then the unmoved mover's role changes. It is not the initiator of a process. It is the necessary ground of an eternal one. The wall at the foundation of The Medium is precisely here: the existence of the eternal generative ground is not derivable within the system it generates. This is acknowledged openly. It is the Gödelian limit of the framework — not a flaw, but the honest boundary of what formal reasoning can reach.
\subsection{What The Medium Is}
\label{sec:org45afad6}
The Medium is not a theory of everything. It is not a competitor to loop quantum gravity or causal set theory on technical grounds. It is a foundational framework that starts from the most primitive possible assumption, derives upward without inserting unjustified structure, locates its wall precisely and honestly, and asks only that what follows from the primitive be consistent with what is observed.
Its primitive is \(\iota\) — a unit occurrence with binary sign disposition. From this single undefined term, and a single primitive relation of co-manifestation, the framework derives the structures that constitute physical reality at the substrate level. The Hamiltonian is derived, not assumed. Stable particle configurations emerge as local extrema of total configuration energy. The projection layer — where dimensional quantities and observable physics live — is explicitly distinguished from the substrate.
Every assumption is stated. Every derivation is shown. The wall is located precisely rather than hidden inside technical language. This is the discipline the framework demands of itself — and the standard by which it asks to be evaluated.
The Medium offers a different position. From that position, different paths are visible. Whether those paths lead to solutions that have remained inaccessible from existing positions is a question the framework opens rather than answers. That is what a new starting point does. It does not arrive. It begins.
\end{document}
